\chapter*{Dedica delle offerte}

\delegateSetUseNext

%\suttaref{Trad.}%
[Yo so] bha꜕gavā a꜕rahaṁ sammāsambuddho\\

\begin{english}
[Al Beato] al Nobile Perfettamente Risvegliato,
\end{english}

Svākkhā꜓to yena bha꜕gava꜓tā dhammo\\

\begin{english}
Alla Verità che ha così ben esposto,
\end{english}


Supaṭi꜕panno yassa bha꜕gava꜕to sāvaka꜕saṅgho\\

\begin{english}
Ai discepoli del Beato, che hanno ben praticato,
\end{english}


Tam-ma꜓yaṁ bha꜕gavantaṁ sa꜕dhammaṁ sa꜕saṅghaṁ\\

\begin{english}
Al Buddha, al Dhamma e al Sangha,
\end{english}


Imehi꜓ sakkārehi꜕ yathārahaṁ āropi꜕tehi a꜕bhi꜓pūja꜕yāma\\

\begin{english}
Offriamo i nostri onori e i nostri omaggi.
\end{english}


Sādhu꜓ no bhante bha꜕gavā su꜕cira-parinibbu꜕topi\\

\begin{english}
è bene per noi che il Beato, pur avendo ottenuto la liberazione,
\end{english}


Pacchi꜓mā-ja꜕na꜓tānu꜓kampa꜕-mānasā\\

\begin{english}
abbia avuto compassione, delle generazioni future
\end{english}


Ime sakkāre dugga꜕ta꜕-paṇṇākāra꜓-bhūte pa꜕ṭiggaṇhātu\\

\begin{english}
questi umili onori e doni siano accolti, a nostro duraturo beneficio
\end{english}


Amhā꜓kaṁ dīgha꜕rattaṁ hi꜕tāya su꜕khāya\\

\begin{english}
e per la felicità, che questo ci dona.
\end{english}

\clearpage

Arahaṁ sammāsambuddho bha꜕gavā\\

\begin{english}
Il Beato, il Nobile Perfettamente Risvegliato
\end{english}


Buddhaṁ bha꜕gavantaṁ a꜕bhi꜓vādemi 

\begin{english}
rendo omaggio al Buddha, al Beato. 
\instr{Inchino}
\end{english}



[Svākkhā꜓to] bha꜕gava꜓tā dhammo\\
Dhammaṁ namassāmi 

\begin{english}
[L'Insegnamento] completamente esposto dal Beato,\\
mi inchino al Dhamma.
  \instr{Inchino}
\end{english}


[Supaṭi꜕panno] bha꜕gava꜕to sāvaka꜕saṅgho\\
Sa꜓ṅghaṁ na꜕māmi 

\begin{english}
 [I discepoli del Beato] che hanno ben praticato,\\
mi inchino al Sangha.
  \instr{Inchino}
\end{english}

\clearpage


\chapter*{Omaggio preliminare}

\begin{leader}
  [Ha꜓nda mayaṁ buddhassa꜕ bhagavato pubbabhāga-namakā꜕raṁ karomase]
\end{leader}

\begin{english}
[Rendiamo ora omaggio preliminare al Buddha]
\end{english}

%\suttaref{DN 21}%
Namo tassa bha꜕gava꜕to araha꜕to sa꜓mmāsa꜓mbuddha꜕ssa
\begin{english}
Omaggio al Beato, al Nobile Perfettamente Risvegliato
\end{english}
\instr{Tre volte}

\chapter*{Riflessione sul Buddha}

\delegateSetUseNext

\begin{leader}
  [Ha꜓nda mayaṁ buddhānu꜕ssa꜕ti꜕nayaṁ karomase]
\end{leader}

\begin{english}
[Ora ricordiamo le qualità del Buddha]
\end{english}

%\suttaref{DN 2}%
Taṁ kho꜓ pana bha꜕gavantaṁ evaṁ kaly꜓āṇo kitti꜕saddo abbhugga꜕to\\
\begin{english}
Così è stato detto del Beato, con grande ammirazione:
\end{english}
Itipi so bha꜕gavā a꜕rahaṁ sammāsambuddho\\
\begin{english}
puro di cuore, Perfettamente Risvegliato e Beato,
\end{english}
Vijjāca꜕raṇa꜓-sampanno su꜕ga꜕to loka꜕vi꜓dū\\
\begin{english}
impeccabile in condotta e conoscenza, il Sereno, Conoscitore dei mondi,
\end{english}
Anu꜓tta꜕ro purisa꜕damma-sārathi satthā deva-ma꜕nussānaṁ\\
\begin{english}
Abile insegnante per chi cerca educazione, Maestro di dèi e di uomini;
\end{english}
\vin buddho bha꜕gavā'ti
\begin{english}
Risvegliato e Beato.
\end{english}
\clearpage


\chapter*{Lode del Buddha}

\delegateSetUseNext

\begin{leader}
  [Ha꜓nda mayaṁ buddhābhi꜕gī꜕tiṁ karomase]
\end{leader}

\begin{english}
[Cantiamo ora la lode del Buddha]
\end{english}

%\suttaref{Trad.}%
Buddh'vāra꜕ha꜓nta-varatādi꜕gu꜓ṇābhi꜕yutto\\

\begin{english}
Il Buddha, dotato di eccellenti qualità\\
\end{english}

Suddhābhi꜕ñāṇa-ka꜕ru꜓ṇāhi sa꜓māga꜕tatto\\

\begin{english}
La sua natura è purezza, conoscenza diretta e compassione\\
\end{english}

Bodhesi꜕ yo su꜕ja꜕na꜓taṁ ka꜕ma꜓laṁ va꜕ sūro\\

\begin{english}
Il Buddha sveglia i cuori come il sole risplende sul loto \\
\end{english}

Vandām'aha꜓ṁ ta꜕m-ara꜕ṇaṁ si꜕rasā꜓ ji꜕nendaṁ\\
\begin{english}
Saluto rispettoso, quel conquistatore di pace\\
\end{english}

Buddho yo sabba꜕-pāṇīnaṁ sa꜕raṇaṁ khema꜕m-utta꜕maṁ\\

\begin{english}
Il Buddha è per tutti i viventi rifugio sicuro e supremo\\
\end{english}

Pa꜕ṭhamānussa꜕tiṭṭhānaṁ vandāmi꜕ taṁ si꜓ren'a꜕haṁ\\

\begin{english}
A questo primo oggetto di ricordo, offro la mia profonda devozione\\
\end{english}

Buddhassā꜓h'a꜕smi dāso/dāsī va buddho me sā꜕mi-ki꜓ssaro\\

\begin{english}
Sono al servizio del Buddha, il Buddha mi ispira e mi guida\\
\end{english}

Buddho dukkhassa꜕ ghātā ca꜕ vidhātā ca꜕ hi꜓tassa꜕ me\\

\begin{english}
Il Buddha mette fine al dolore ed è fonte di ogni mia gioia\\
\end{english}

Buddhass'āha꜓ṁ niyyādemi sa꜕rīrañ-jīvi꜕tañ-ci꜕daṁ\\

\begin{english}
Al Buddha offro il mio corpo e la mia vita\\
\end{english}

Vandanto'ha꜓ṁ/Vandantī'ha꜓ṁ ca꜕rissāmi buddhass'eva꜕ su꜓bodhi꜕taṁ\\

\begin{english}
E con devozione camminerò la sua via di Risveglio\\
\end{english}

Natthi me sa꜕ra꜓ṇaṁ aññaṁ buddho me sa꜕ra꜓ṇaṁ va꜕raṁ\\

\begin{english}
Per me non c'è altro rifugio, il Buddha è il miglior rifugio\\
\end{english}

Etena sacca꜕-vajjena vaḍḍheyyaṁ sa꜕tthu-sā꜓sane\\

\begin{english}
In forza della mia sincerità, che io cresca nella via del maestro\\
\end{english}

Buddhaṁ me vanda꜕mānena/vanda꜕mānāya
yaṁ puññaṁ pa꜕su꜓taṁ i꜕dha\\

\begin{english}
 possa il bene generato da questa lode del Buddha \\
\end{english}

Sa꜕bbepi anta꜕rāyā me māhe꜓su꜓ṁ ta꜕ssa꜓ teja꜕sā

\begin{english}
  aiutarmi a superare ogni difficoltà.
\end{english}

\begin{instruction}
  STANDO CHINATO
\end{instruction}

Kāyena vācāya va ceta꜕sā꜓ vā\\
\begin{english}
se con il corpo, la parola o la mente \\
\end{english}

Bu꜓ddhe ku꜕kammaṁ pa꜕kataṁ ma꜕yā yaṁ\\
\begin{english}
ho agito male nei confronti del Buddha\\
\end{english}

Bu꜓ddho pa꜕ṭiggaṇhā꜕tu acca꜕yantaṁ\\

\begin{english}
si accetti il riconoscimento del mio errore \\
\end{english}

Kālantare sa꜓ṁvarituṁ va꜕ bu꜓ddhe

\begin{english}
  e l'intenzione di correggermi nei riguardi del Buddha.
\end{english}

\clearpage


\chapter*{Riflessione sul Dhamma}

\delegateSetUseNext

\begin{leader}
  [Ha꜓nda mayaṁ dhammānu꜕ssa꜕ti꜕nayaṁ karomase]
\end{leader}

\begin{english}
[Ora ricordiamo le qualità del Dhamma]
\end{english}

%\suttaref{DN 16}%
Svākkhā꜓to bha꜕gava꜓tā dhammo\\
\begin{english}
Il Dhamma è ben esposto dal Beato\\
\end{english}

Sa꜓ndiṭṭhi꜕ko a꜕kāli꜕ko ehi꜕passi꜕ko\\
\begin{english}
Presente qui ed ora, senza-tempo, incoraggia la verifica diretta\\
\end{english}


Opanayi꜕ko pa꜕cca꜕ttaṁ vedi꜓ta꜕bbo viññūhī'ti
\begin{english}
Conduce oltre, va vissuto di persona da ciascun ricercatore.\\
\end{english}

\chapter*{Lode del Dhamma}

\begin{leader}
  [Ha꜓nda mayaṁ dhammābhi꜕gī꜕tiṁ karomase]
\end{leader}

\begin{english}
[Cantiamo ora la lode del Dhamma]
\end{english}


%\suttaref{Trad.}%
Svākkhā꜓ta꜕t'ādi꜕guṇa-yoga꜕-va꜓sena꜕ seyyo\\

\begin{english}
È eccellente perché è ben esposto \\
\end{english}

Yo magga꜕-pāka-pa꜕riyatti꜕-vi꜓mokkha꜕-bhedo\\

\begin{english}
Nei suoi diversi aspetti: Sentiero e Frutto, pratica e liberazione.\\
\end{english}

Dhammo ku꜕loka-pa꜕ta꜓nā ta꜕da꜓-dhāri꜕-dhārī\\

\begin{english}
Coloro che lo sostengono, sono a loro volta protetti dal cadere nell'illusione.\\
\end{english}

Vandām'aha꜓ṁ ta꜕ma-ha꜕raṁ va꜕ra-dha꜓mma꜕m-etaṁ\\

\begin{english}
Mi inchino alla Verità preziosa che disperde l'ignoranza.\\
\end{english}

Dhammo yo sabba꜕-pāṇīnaṁ sa꜕raṇaṁ khema꜕m-utta꜕maṁ\\

\begin{english}
Il Dhamma è per tutti i viventi rifugio sicuro e supremo.\\
\end{english}

Du꜕tiyānussa꜕tiṭṭhānaṁ vandāmi꜕ taṁ si꜓ren'a꜕haṁ\\

\begin{english}
A questo secondo oggetto di ricordo offro la mia profonda devozione\\
\end{english}

Dhammassā꜓h'a꜕smi dāso/dāsī va dhammo me sā꜕mi-ki꜓ssaro\\

\begin{english}
Sono al servizio del Dhamma, il Dhamma mi ispira e mi guida\\
\end{english}

Dhammo dukkhassa꜕ ghātā ca꜕ vidhātā ca꜕ hi꜓tassa꜕ me\\

\begin{english}
Il Dhamma mette fine al dolore ed è fonte di ogni mia gioia\\
\end{english}

Dhammass'āha꜓ṁ niyyādemi sa꜕rīrañ-jīvi꜕tañ-ci꜕daṁ\\

\begin{english}
Al Dhamma offro il mio corpo e la mia vita\\
\end{english}

Vandantoha꜓ṁ/Vandantīha꜓ṁ ca꜕rissāmi dhammass'eva꜕ su꜓dhamma꜕taṁ\\

\begin{english}
E con devozione camminerò l'eccellente via del Dhamma\\
\end{english}

Natthi me sa꜕ra꜓ṇaṁ aññaṁ dhammo me sa꜕ra꜓ṇaṁ va꜕raṁ\\

\begin{english}
Per me non c'è altro rifugio, il Dhamma è il miglior rifugio\\
\end{english}

Etena sacca꜕-vajjena vaḍḍheyyaṁ sa꜕tthu-sā꜓sane\\

\begin{english}
In forza della mia sincerità che io cresca nella via del Maestro\\
\end{english}

% ...existing code...
\mbox{Dhammaṁ me vanda꜕mānena/vanda꜕mānāya yaṁ puññaṁ pa꜕su꜓taṁ i꜕dha}\\
% ...existing code...

\begin{english}
 possa il bene generato da questa lode del Dhamma \\
\end{english}


Sa꜕bbepi anta꜕rāyā me māhe꜓su꜓ṁ ta꜕ssa꜓ teja꜕sā
\begin{english}
  aiutarmi a superare ogni difficoltà.
\end{english}



\clearpage


\begin{instruction}
  STANDO CHINATO
\end{instruction}

Kāyena vācāya va ceta꜕sā꜓ vā\\
\begin{english}
se con il corpo, la parola o la mente \\
\end{english}

Dha꜓mme ku꜕kammaṁ pa꜕kataṁ ma꜕yā yaṁ\\
\begin{english} 
ho agito male nei confronti del Dhamma\\
\end{english} 

Dha꜓mmo pa꜕ṭiggaṇhā꜕tu acca꜕yantaṁ\\
\begin{english}
si accetti il riconoscimento del mio errore \\
\end{english} 

Kālantare sa꜓ṁvarituṁ va꜕ dha꜓mme
\begin{english}
  e l'intenzione di correggermi nei riguardi del Dhamma.
\end{english}

\clearpage


\chapter*{Riflessione sul Saṅgha}

\delegateSetUseNext

\begin{leader}
  [Ha꜓nda mayaṁ saṅghānu꜕ssa꜕ti꜕nayaṁ karomase]
\end{leader}

\begin{english}
[Ora ricordiamo le qualità del Sangha]
\end{english}

%\suttaref{DN 16}%
Supaṭi꜕panno bha꜕gava꜕to sāvaka꜕saṅgho\\

\begin{english}
I discepoli del Beato che hanno ben praticato\\
\end{english}

Ujupaṭi꜕panno bha꜕gava꜕to sāvaka꜕saṅgho\\
\begin{english}
Vivono rettamente\\
\end{english}

Ñāyapaṭi꜕panno bha꜕gava꜕to sāvaka꜕saṅgho\\

\begin{english}
Vivono sapientemente\\
\end{english}

Sā꜓mīci꜕pa꜕ṭi꜕panno bha꜕gava꜕to sāvaka꜕saṅgho\\

\begin{english}
Vivono con integrità\\
\end{english}



Yadidaṁ cattāri purisa꜕yugāni aṭṭha꜓ purisa꜕pugga꜕lā\\
\begin{english}
Sono le quattro coppie, le otto specie di Nobili persone\\
\end{english}



Esa bha꜕gava꜕to sāvaka꜕saṅgho\\
\begin{english}
Tali sono i discepoli del Beato\\
\end{english}



Āhu꜕neyyo pāhu꜕neyyo dakkhi꜕ṇeyyo añja꜕li-ka꜕ra꜓ṇīyo\\
\begin{english}   
Sono degni di doni, degni di ospitalità, degni di offerte, degni di rispetto\\
\end{english}


Anu꜓tta꜕raṁ puññakkhe꜕ttaṁ lokassā'ti
\begin{english}
Il più fertile terreno di meriti per il mondo.
\end{english}

\chapter*{Lode del Saṅgha}

\begin{leader}
  [Ha꜓nda mayaṁ saṅghābhi꜕gī꜕tiṁ karomase]
\end{leader}

\begin{english}
[Cantiamo ora la lode del Saṅgha]
\end{english}

%\suttaref{Trad.}%
Sa꜕ddhammajo supaṭipatti꜕-gu꜓ṇādi꜕yutto\\
\begin{english}
Nato dal Dhamma, il Saṅgha che ha praticato bene\\
\end{english}

Yo'ṭṭhabbi꜕dho ari꜓yapugga꜕la꜓-saṅgha꜕-seṭṭho\\
\begin{english}
Quel Saṅgha che è composto da otto specie di Nobili persone.\\
\end{english}


Sī꜓lādi꜕dhamma-pa꜕varāsa꜕ya꜓-kāya꜕-citto\\
\begin{english}
Corpo e mente, guidati da saggezza e integrità morale.\\
\end{english}


Vandām'aha꜓ṁ ta꜕m-ari꜕yāna꜕-gaṇa꜓ṁ su꜕suddhaṁ\\
\begin{english}
Rendo omaggio alla purezza di questa Nobile assemblea.\\
\end{english}   

Sa꜓ṅgho yo sabba꜕-pāṇīnaṁ sa꜕raṇaṁ khema꜕m-utta꜕maṁ\\
\begin{english}
Il Sangha è per tutti i viventi rifugio sicuro e supremo.\\
\end{english}

Ta꜕tiyānussa꜕tiṭṭhānaṁ vandāmi꜕ taṁ si꜓ren'a꜕haṁ
\begin{english}
A questo terzo oggetto di ricordo offro la mia profonda devozione\\
\end{english}

Saṅghass'ā꜓ha꜕smi dāso/dāsī va saṅgho me sā꜕mi-ki꜓ssaro\\
\begin{english}
Sono al servizio del Sangha, il Sangha mi ispira e mi guida\\
\end{english} 

Sa꜓ṅgho dukkhassa꜕ ghātā ca꜕ vi꜓dhātā ca꜕ hi꜓tassa꜕ me\\
\begin{english}
Il Sangha mette fine al dolore ed è fonte di ogni mia gioia\\
\end{english}

Saṅghass'āha꜓ṁ niyyādemi sa꜕rīrañ-jīvi꜕tañ-ci꜕daṁ\\
\begin{english}
Al Sangha offro il mio corpo e la mia vita\\
\end{english}

Vandanto'ha꜓ṁ/Vandantī'ha꜓ṁ ca꜕rissāmi saṅghassopa꜕ṭi꜓panna꜕taṁ\\
\begin{english}
E con devozione camminerò la corretta condotta del Sangha\\
\end{english}

Natthi me sa꜕ra꜓ṇaṁ aññaṁ saṅgho me sa꜕ra꜓ṇaṁ va꜕raṁ\\
\begin{english}
Per me non c'è altro rifugio, il Sangha è il miglior rifugio\\
\end{english}

Etena sacca꜕-vajjena vaḍḍheyyaṁ sa꜕tthu-sā꜓sane\\
\begin{english}
In forza della mia sincerità che io cresca nella via del Maestro\\
\end{english}

Sa꜓ṅghaṁ me vanda꜕mānena/vanda꜕mānāya\\
\vin yaṁ puññaṁ pa꜕su꜓taṁ i꜕dha\\
\begin{english}
 possa il bene generato da questa lode del Sangha \\
\end{english}


Sa꜕bbepi anta꜕rāyā me māhe꜓su꜓ṁ ta꜕ssa꜓ teja꜕sā
\begin{english}
  aiutarmi a superare ogni difficoltà.
\end{english}

\begin{instruction}
  STANDO CHINATO
\end{instruction}

Kāyena vācāya va ceta꜕sā꜓ vā\\
\begin{english}
se con il corpo, la parola o la mente \\
\end{english}
Sa꜓ṅghe ku꜕kammaṁ pa꜕kataṁ ma꜕yā yaṁ\\
\begin{english}
ho agito male nei confronti del Sangha\\
\end{english}
Sa꜓ṅgho pa꜕ṭiggaṇhā꜕tu acca꜕yantaṁ\\
\begin{english}
si accetti il riconoscimento del mio errore \\
\end{english}
Kālantare sa꜓ṁvarituṁ va꜕ sa꜓ṅghe
\begin{english}
  e l'intenzione di correggermi nei riguardi del Sangha.
\end{english}


\begin{instruction}
  A questo punto si pratica la meditazione in silenzio, a volte seguita da un discorso di Dhamma, e si conclude con quanto segue:
\end{instruction}

\chapter*{Omaggio finale}

\delegateSetUseNext

%\suttaref{Trad.}%
[Arahaṁ] sammāsambuddho bha꜕gavā\\
Buddhaṁ bha꜕gavantaṁ a꜕bhi꜓vādemi

\begin{english}
[Il Beato] il Nobile Perfettamente Risvegliato
rendo omaggio al Buddha, al beato\\ 
\end{english}
\instr{Inchino}

[Svākkhā꜓to] bha꜕gava꜓tā dhammo\\
Dhammaṁ namassāmi 
\begin{english}
[L'Insegnamento] completamente esposto dal Beato\\
mi inchino al Dhamma
\end{english}
\instr{Inchino}

[Supaṭi꜕panno] bha꜕gava꜕to sāvaka꜕saṅgho\\
Sa꜓ṅghaṁ na꜕māmi  
\begin{english}
 [I discepoli del Beato] che hanno ben praticato\\
mi inchino al Sangha
\end{english}
\instr{Inchino}

\clearpage

% End of evening-chanting.tex
